\section{Отзыв сертификата}\label{ACME.Revoke}

Для отзыва сертификата клиент отправляет POST-запрос на узел \code{revokeCert}. 
Тело запроса представляет собой подписанный объект JSON со следующими полями:

\begin{itemize}
\item
\sstr{certificate} (обязательное, строка)~--- 
сертификат, который необходимо отозвать. Указывается DER-код сертификата,
повторно закодированный по правилам base64url;

\item
\sstr{reason} (необязательно, целое)~--- код причины отзыва в соответствии 
с СТБ~34.101.19 (п. 7.3.1, расширение \code{reasonCode}).
%
Если это поле не задано, УЦ, агентом которого является сервер ACME, должен 
опускать расширение \code{reasonCode} при формировании ответов OCSP и СОС. 

Сервер может запретить клиенту использовать определенные коды причин отзыва.
Если запрос содержит запрещенный код, то сервер должен отклонить запрос, указав 
тип ошибки \sstr{badRevocationReason}. В подробностях проблемы следует
указать, какие коды разрешены. 
\end{itemize}

Запросы на отзыв отличаются от других запросов ACME тем, что они могут быть 
подписаны либо \pdfexample{ключом учетной записи}{acme/revoke1.txt}, либо 
\pdfexample{ключом сертификата}{acme/revoke2.txt}.  
%
В первом случае в заголовке объекта JWS открытый ключ описывается через 
поле \sstr{kid}, во втором~--- через поле \sstr{jwk}.

\if0
Пример использования пары ключей учетной записи для подписи: 

\begin{codeblock}
POST /acme/revoke-cert HTTP/1.1
Host: example.com
Content-Type: application/jose+json

{
  "protected": base64url({
    "alg": "BIGNS128",
    "kid": "https://example.com/acme/acct/evOfKhNU60wg",
    "nonce": "JHb54aT_KTXBWQOzGYkt9A",
    "url": "https://example.com/acme/revoke-cert"
  }),
  "payload": base64url({
    "certificate": "MIIEDTCCAvegAwIBAgIRAP8...",
    "reason": 4
  }),
  "signature": "Q1bURgJoEslbD1c5...3pYdSMLio57mQNN4"
}
\end{codeblock}

Пример использования пары ключей сертификата для подписи:

\begin{codeblock}
POST /acme/revoke-cert HTTP/1.1
Host: example.com
Content-Type: application/jose+json

{
  "protected": base64url({
    "alg": "BIGNS128",
    "jwk": {...},
    "nonce": "JHb54aT_KTXBWQOzGYkt9A",
    "url": "https://example.com/acme/revoke-cert"
  }),
  "payload": base64url({
    "certificate": "MIIEDTCCAvegAwIBAgIRAP8...",
    "reason": 1
  }),
  "signature": "Q1bURgJoEslbD1c5...3pYdSMLio57mQNN4"
} 
\end{codeblock}
\fi

Перед отзывом сертификата сервер должен убедиться, что личный ключ,
использованный для подписи запроса на отзыв, авторизован для отзыва. Сервер
должен принимать в качестве авторизованных, по крайней мере, следующие ключи:

\begin{itemize}
\item 
личный ключ учетной записи, по запросу владельца которой выпущен сертификат;
\item
личный ключ учетной записи, которая была авторизована относительно всех 
идентификаторов сертификата;
\item
личный ключ, соответствующий открытому ключу сертификата.
\end{itemize}

Если отзыв выполнен, сервер возвращает \pdfexample{ответ}{acme/revoke3.txt} с 
кодом статуса 200. Если отзыв не состоялся, сервер возвращает 
\pdfexample{ошибку}{acme/revoke4.txt}. Например, если сертификат уже был 
отозван, сервер возвращает ответ с кодом статуса 400 (неверный запрос) и типом 
\sstr{alreadyRevoked}.

\if0
\begin{codeblock}
HTTP/1.1 200 OK
Replay-Nonce: IXVHDyxIRGcTE0VSblhPzw
Content-Length: 0
Link: <https://example.com/acme/directory>;rel="index"
\begin{codeblock}

\begin{codeblock}
HTTP/1.1 403 Forbidden
Replay-Nonce: lXfyFzi6238tfPQRwgfmPU
Content-Type: application/problem+json
Content-Language: en
Link: <https://example.com/acme/directory>;rel="index"

{
  "type": "urn:ietf:params:acme:error:unauthorized",
  "detail": "No authorization provided for name example.org"
}
\end{codeblock}
\fi
